\chapter{Einführung}
\label{sec:einführung}

Das menschliche Fahrverhalten ist bedingt durch die Vielzahl der möglichen Situationen im Straßenverkehr vielseitig. Um das Verhalten in verschiedenen Situationen zu untersuchen, müssen entweder Daten aus dem realen Verkehr ausgewertet werden oder bestimmte Situationen in einer kontrollierten Testumgebung nachgestellt werden. In einer realen Testumgebung ist es allerdings schwer, bestimmte Situationen gezielt nachzustellen. Die Umgebungsbedingungen des realen Verkehrs mit unterschiedlichen Fahrumgebungen wie Landstraßen, Autobahnen, Städten oder verschiedene Verkehrsteilnehmern wie Autofahrern, Straßenbahnen, Radfahrern und Fußgängern sind komplex. Ein anderer Ansatz ist es, die Umgebung zu simulieren. Auch für die Nachbildung von gefährlichen Situationen im Straßenverkehr weist eine Simulation aus finanzieller und praktischer Sicht Vorteile auf. Verschiedene Verkehrssituationen können mit kleineren finanziellen Aufwänden wiederholbar und zeitlich flexibel durchgeführt werden. Das Automotive Mobility Lab des Institus für energieeffiziente Mobilität ({\acs{IEEM}}) entwickelt aus diesem Grund ein Simulationsframework, welches mit unterschiedlichen Fortbewegungsmethoden (Auto, Fahrrad, Gehen) und in unterschiedlicher Realitätsnähe genutzt werden kann. Dadurch sind \Gls{Human-in-the-Loop}-Fahrversuche möglich, bei denen Fahrversuche mit einer realen Testperson in einer simulierten Umgebung durchgeführt werden können.  

Als Softwarebasis dient das Open-Source-Projekt \Gls{CARLA}. Dieses ermöglicht es, verschiedene Verkehrsszenarien nachzustellen. Die Analyse des menschlichen Verhaltens im Straßenverkehr kann somit ohne Gefahren und in vielfältigen simulierten Szenarien erfolgen. Die \Gls{Server-Client Struktur} von CARLA lässt ebenso mehrere Probanden in demselben Verkehrsszenario zu. So können menschliche Interaktionen realitätsnah nachgebildet werden.

Um die Ursachen bestimmter Entscheidungen im Straßenverkehr zu ergründen, soll in dem Fahrsimulator ein Eye Tracking System integriert werden. Durch die Auswertung der Blickposition des Fahrers sollen Informationen zum Wahrnehmungs- und Entscheidungsprozess des Fahrers in gewissen Situationen gewonnen werden. Hierdurch soll deutlich werden, was die Probanden während der Testfahrt sehen, worauf diese sich konzentrieren und welche Handlungsoptionen sie in Gefahrensituationen abwägen.

\newpage

\section{Zielsetzung}

Im Rahmen dieser Projektarbeit soll ein kamera basierter Eye Tracker in den Fahrsimulator der Hochschule bzw. in die CARLA Simulationsumgebung integriert werden.

Dazu sollen verschiedene verfügbare Eye Tracker auf die Anwendbarkeit und Integrierbarkeit in CARLA bewertet werden. Zur Erprobung des Eye Trackers soll ein Testszenario in CARLA erstellt werden. Weitere Vorgaben sind die Projektion des Blickpunktes auf die Leinwand und eine Ausgabe der aufgezeichneten Daten. Außerdem sollen die Daten innerhalb der Simulationsumgebung mit den Daten des Eye Trackers vereint werden. Dafür soll vor allem der \Gls{Instance Segmentation} Sensor innerhalb von CARLA zum Einsatz kommen. Dieser wird in einem späteren Abschnitt genauer erklärt. Ziel wird es sein, die Daten des Eye Trackers mit den Daten dieses Sensors zu vereinen, um Aussagen über die von der Testperson angeschauten Objekte in der Simulation zu erhalten.  


\section{Lösungsansatz}

Zur Zielerreichung müssen zuerst die Grundlagen von Eye Tracker Systemen und der CARLA-Simulationsumgebung erarbeitet werden. Für die Integration in das Gesamtsystem soll ein objektorientierter Python Code entwickelt werden. Damit die Ziele erreicht werden können, ist vor allem eine intensive Auseinandersetzung mit den Application Programming Interfaces ({\acs{API}})s der verschiedenen beteiligten Systemen notwendig. Eine Eigenentwicklung eines Eye Tracking Systems ist technisch und mathematisch anspruchsvoll und würde den zeitlichen Rahmen einer Projektarbeit übersteigen. Aus diesem Grund muss ein bereits vorhandenes System gefunden werden, welches kostengünstig ist und welches mindestens eine {\acs{API}} bietet, mit welcher die ausgegebenen Daten weiterverarbeitet werden können. Außerdem muss sichergestellt werden, dass die Krümmung der Leinwand bei der Verwendung des Eye Tracking Systems nicht zu Ungenauigkeiten führt. Entweder wird ein Eye Tracking System eingesetzt, welches gekrümmte Leinwände nativ unterstützt oder es muss eine selbst entwickelte Lösung erarbeitet werden. 






